% P11 paper module: R20 individual noncommutativity versus relative polar incompatibility.

\subsection{Individual metric noncommutativity versus relative polar incompatibility}
\label{subsec:o3s-relative-polar}

Keep the odd-sector notation of Subsections~\ref{subsec:o3n-gauge-intertwining} and
\ref{subsec:o3r-sqrt-offblock}.  For a single source level $X\in\{R,S\}$ write
\[
 B_X:=G_{X,T_0}^-,\qquad C_X:=G_{X,U}^-,
\]
\[
 X_X=C_X^{1/2}B_X^{-1/2}=U_XH_X,
 \qquad
 H_X:=\bigl(B_X^{-1/2}C_XB_X^{-1/2}\bigr)^{1/2}.
\]
Proposition~\ref{prop:o3n-gauge-commutator} gives the qualitative equivalence
$U_X=I\iff[B_X,C_X]=0$.  The first point below is the scale-correct quantitative version
for one metric pair.

\begin{proposition}[Normalized square-root commutator controls individual polar activity]
\label{prop:o3s-individual-commutator}
For any boundedly invertible positive pair $B,C$, let
\[
 X=C^{1/2}B^{-1/2}=UH,
 \qquad
 H=(B^{-1/2}CB^{-1/2})^{1/2}.
\]
Then
\[
 X-X^*
 =B^{-1/2}[B^{1/2},C^{1/2}]B^{-1/2}
\tag{O3S.1}\label{eq:o3s-skew}
\]
and
\[
 \boxed{
 \|U-I\|
 \ge
 \frac{
 \|B^{-1/2}[B^{1/2},C^{1/2}]B^{-1/2}\|
 }{2\|H\|}.
 }
\tag{O3S.2}\label{eq:o3s-U-lower}
\]
The right-hand side is invariant under the common rescaling $(B,C)\mapsto(tB,tC)$,
$t>0$.
\end{proposition}

\begin{proof}
The identity~\eqref{eq:o3s-skew} is immediate from
\[
 B^{-1/2}[B^{1/2},C^{1/2}]B^{-1/2}
 =C^{1/2}B^{-1/2}-B^{-1/2}C^{1/2}.
\]
Since $X=UH$ and $X^*=HU^*$,
\[
 X-X^*=(U-I)H+H(I-U^*).
\]
Hence
\[
 \|X-X^*\|\le2\|H\|\,\|U-I\|,
\]
which proves~\eqref{eq:o3s-U-lower}.  Under common scalar rescaling both $X$ and $H$
are unchanged, and so is the normalized expression.
\end{proof}

The actual P11 gauge is not the individual quantity $U_X-I$ but the relative defect
\[
 \mathfrak P_U:=U_SW-WU_R.
\tag{O3S.3}\label{eq:o3s-relative-defect}
\]
Let $P=WW^*$ and retain
\[
 \Gamma_U=W^*U_SWU_R^*.
\]

\begin{proposition}[Exact relative polar decomposition]
\label{prop:o3s-relative-decomposition}
Put
\[
 \Lambda_U:=(I-P)U_SW.
\]
Then
\[
 \boxed{
 \mathfrak P_U
 =\Lambda_U+W(\Gamma_U-I)U_R,
 }
\tag{O3S.4}\label{eq:o3s-relative-decomp}
\]
and the two summands have orthogonal ranges.  Consequently, for every source vector
$f$,
\[
 \boxed{
 \|\mathfrak P_Uf\|^2
 =\|\Lambda_Uf\|^2
  +\|(\Gamma_U-I)U_Rf\|^2.
 }
\tag{O3S.5}\label{eq:o3s-relative-pythagoras}
\]
In particular a lower bound for the polar leakage $\Lambda_U$ would be sufficient for a
lower bound on the true relative gauge defect, but a lower bound on either individual
$\|U_R-I\|$ or $\|U_S-I\|$ is not.
\end{proposition}

\begin{proof}
Using $P=WW^*$ and the definition of $\Gamma_U$,
\[
\begin{aligned}
 \mathfrak P_U
 &=(I-P)U_SW+PU_SW-WU_R\\
 &=\Lambda_U+W(W^*U_SW-U_R)\\
 &=\Lambda_U+W(\Gamma_U-I)U_R.
\end{aligned}
\]
The first term lies in $(\Ran W)^\perp$ and the second in $\Ran W$, proving the
orthogonality and~\eqref{eq:o3s-relative-pythagoras}.
\end{proof}

The next model shows that individual metric noncommutativity, even if arbitrarily large
and even if it produces genuinely nontrivial polar factors, does not force relative
polar incompatibility.

\begin{proposition}[Canonical nested noncommutativity can cancel perfectly]
\label{prop:o3s-matched-countermodel}
There is a canonical inclusion model with exact base/future metric pullback in which
both source and target metric pairs are noncommuting, their raw commutator norms can be
made arbitrarily large, and $U_R,U_S$ are nontrivial, but
\[
 \boxed{
 (I-P)U_SW=0,
 \qquad
 U_SW-WU_R=0
 }
\tag{O3S.6}\label{eq:o3s-matched-zero}
\]
identically.
\end{proposition}

\begin{proof}
Let
\[
 B=\begin{pmatrix}1&0\\0&2\end{pmatrix},
 \qquad
 C=\begin{pmatrix}2&1\\1&2\end{pmatrix}.
\]
Then $B,C>0$ and
\[
 [B,C]=\begin{pmatrix}0&-1\\1&0\end{pmatrix}\ne0.
\]
Take source $\mathbb C^2$, target $\mathbb C^3$, and the canonical inclusion
$Wz=(z,0)$.  For $t>0$ set
\[
 G_{R,T_0}=tB,
 \qquad G_{R,U}=tC,
\]
\[
 G_{S,T_0}=tB\oplus1,
 \qquad G_{S,U}=tC\oplus1.
\]
The base and future metric pullback identities hold exactly.  Moreover the normalized
base transition is precisely the canonical inclusion $W$.  The polar products satisfy
\[
 X_S=X_R\oplus1,
\]
so uniqueness of polar decomposition gives
\[
 U_S=U_R\oplus1.
\]
Since $[B,C]\ne0$, Proposition~\ref{prop:o3n-gauge-commutator} gives $U_R\ne I$ and
$U_S\ne I$.  Nevertheless $U_SW=WU_R$, proving~\eqref{eq:o3s-matched-zero}.  Finally
\[
 \|[G_{R,T_0},G_{R,U}]\|
 =t^2\|[B,C]\|,
\]
and the same holds on the target nontrivial block, so the raw commutators are
arbitrarily large.
\end{proof}

The previous model has no modulus leakage.  The following direct-sum strengthening
shows that adding an R19-type modulus defect does not repair the implication.

\begin{proposition}[Modulus leakage plus large individual commutators still need not produce polar leakage]
\label{prop:o3s-combined-countermodel}
There is a canonical inclusion model satisfying exact base/future pullback such that
\[
 \|(I-P)Q\|>0,
\]
its square-root off-block can be made arbitrarily large, and both individual metric-pair
commutators can simultaneously be made arbitrarily large, while nevertheless
\[
 \boxed{
 (I-P)U_SW=0,
 \qquad
 U_SW-WU_R=0.
 }
\tag{O3S.7}\label{eq:o3s-combined-zero}
\]
\end{proposition}

\begin{proof}
For the first direct summand take the R19 canonical inclusion block
\[
 W_1:\mathbb C\hookrightarrow\mathbb C^2,
 \qquad W_1z=(z,0),
\]
with base metrics $1$ and $I_2$, and future metrics
\[
 2s^2,
 \qquad
 s^2A_0,
 \qquad
 A_0=\begin{pmatrix}2&1\\1&2\end{pmatrix}.
\]
Its polar factors are $1$ and $I_2$, while
\[
 Q_1=A_0^{1/2}W_1/\sqrt2
\]
has nonzero off-range component because $e_1$ is not an eigenvector of $A_0$; the
corresponding square-root off-block grows linearly in $s$.

For the second summand use the same noncommuting pair $B,C$ as in
Proposition~\ref{prop:o3s-matched-countermodel}, with identity inclusion between two
copies of $\mathbb C^2$, base metric $tB$ and future metric $tC$ on both copies.  Its
source and target polar factors are the same nontrivial unitary $U_0$, while the raw
commutator norm is $t^2\|[B,C]\|$.

Take the orthogonal direct sums of the two systems.  The total transition is a canonical
coordinate inclusion, the base/future pullback identities remain exact, and
\[
 U_R=1\oplus U_0,
 \qquad
 U_S=I_2\oplus U_0.
\]
Hence~\eqref{eq:o3s-combined-zero} holds.  At the same time the first summand supplies
nonzero modulus leakage and arbitrarily large square-root off-block as $s\to\infty$,
while the second supplies arbitrarily large individual commutators as $t\to\infty$.
\end{proof}

\begin{remark}[R20 firewall: the missing object is relative, not individual]
\label{rem:o3s-firewall}
Proposition~\ref{prop:o3s-individual-commutator} gives a legitimate quantitative route
from a scale-normalized square-root commutator to individual polar activity
$\|U_X-I\|$.  Propositions~\ref{prop:o3s-matched-countermodel} and
\ref{prop:o3s-combined-countermodel} show, however, that neither large raw
noncommutativity nor the combination of nontrivial modulus leakage with large
individual noncommutativity can be promoted by abstract algebra to a lower bound on
\[
 (I-P)U_SW
 \quad\text{or}\quad
 U_SW-WU_R.
\]
Thus the proposed target ``prove the two concrete metric pairs are quantitatively
noncommuting'' is insufficient as stated.  The remaining problem for this
polar-gauge obstruction route is a relative polar-compatibility problem: one must
detect a failure of the source and target polar rotations to intertwine under $W$.

Equivalently, the genuinely decisive quantity for this gauge component remains
\[
 \mathfrak P_U=U_SW-WU_R,
\]
or, by Proposition~\ref{prop:o3n-gauge-criterion}, the true gauge compression
$\Gamma_U-I$.  Any useful commutator criterion for this route must therefore be
relative across the nested metric pairs, not a pair of separate estimates on
$[G_{R,T_0},G_{R,U}]$ and $[G_{S,T_0},G_{S,U}]$.

No conclusion about $\Gamma_U\to I$, $K_{R,S}^{T_0,U}\to I$, strong terminal
transport, Object X, Seal, or RH follows here.
\end{remark}
